\section{Role-Based Access Control (RBAC)}

% --- Slide 63: Contents ---
\begin{frame}{Contents}
    \Large
    \begin{enumerate}
        \color{gray}
        \item Introduction to Database Security Issues
        \item Discretionary Access Control (DAC)
        \item Mandatory Access Control (MAC)
        
        \color{black} \bfseries
        \item Role-Based Access Control (RBAC)
        \item Attribute-Based Access Control (ABAC)
        
        \color{gray} \mdseries
          \item Context-Based Access Control (CBAC)
    \end{enumerate}
\end{frame}


% --- Slide 64: Role-Based Access Control (1) ---
\begin{frame}{Role-Based Access Control (1)}
    \Large
    \begin{itemize}
        \item Role-based access control (RBAC) emerged rapidly in the 1990s as a proven technology for managing and enforcing security in large-scale enterprise-wide systems.
        
        \vspace{0.3cm}
        \item Its basic notion is that \textbf{permissions are associated with roles}, and \textbf{users are assigned to appropriate roles}.
        
        \item Roles can be created using the \textbf{CREATE ROLE} and \textbf{DESTROY ROLE} commands.
        
        \begin{itemize}
            \large \item The \textbf{GRANT} and \textbf{REVOKE} commands discussed under DAC can then be used to assign and revoke privileges from roles.
    \end{itemize}
\end{itemize}
\end{frame}

% --- Slide 65: Role-Based Access Control (2) ---
\begin{frame}{Role-Based Access Control (2)}
    \Large
    \begin{itemize}
        \item RBAC appears to be a \textbf{viable alternative} to traditional discretionary and mandatory access controls; it ensures that only authorized users are given access to certain data or resources.
        
        \vspace{0.3cm}
        \item Many DBMSs have allowed the \textbf{concept of roles}, where privileges can be assigned to roles.
        
        \vspace{0.3cm}
        \item \textbf{Role hierarchy} in RBAC is a natural way of organizing roles to reflect the organization's lines of authority and responsibility.
    \end{itemize}
\end{frame}

% --- Slide 66: Role-Based Access Control (3) ---
\begin{frame}{Role-Based Access Control (3)}
    \Large
    \begin{itemize}
        \item Another important consideration in RBAC systems is the possible \textbf{temporal constraints} that may exist on roles, such as time and duration of role activations, and timed triggering of a role by an activation of another role.
        
        \vspace{0.3cm}
        \item Using an RBAC model is highly desirable goal for addressing the key security requirements of \textbf{Web-based applications}.
        
        \item In contrast, DAC and MAC models lack capabilities needed to support the security requirements emerging enterprises and Web-based applications.
    \end{itemize}
\end{frame}

%\section[Mandatory Access Control (MAC)]{Mandatory \\ Access Control (MAC)}
\section[Attribute-Based Access Control]{Attribute-Based \\ Access Control (ABAC)}

% --- Slide: Attribute-Based Access Control (1) ---
\begin{frame}{Attribute-Based Access Control (ABAC) (1)}
    \Large
    \begin{itemize}
        \item \textbf{Attribute-Based Access Control (ABAC)} is a logical access control model that evaluates \textbf{attributes} (characteristics or properties) rather than roles.
        
        \vspace{0.3cm}
        \item Access rights are granted to users through the use of \textbf{policies} which evaluate:
        \begin{itemize}
            \normalsize
            \setlength{\itemsep}{0.2cm}
            \item \textbf{Subject attributes}: details about the user (e.g., job title, department, clearance level).
            \item \textbf{Object attributes}: details about the data/resource (e.g., document type, classification, creation date).
            \item \textbf{Environment conditions}: context of the access request (e.g., time of day, IP address, location).
        \end{itemize}
    \end{itemize}
\end{frame}

% --- Slide: Attribute-Based Access Control (2) - ABAC vs RBAC ---
\begin{frame}{Attribute-Based Access Control (ABAC) (2)}
    \Large
    \begin{itemize}
        \item \textbf{RBAC} uses predefined \textbf{roles}. It can become difficult to manage when the number of roles explodes ("Role Explosion").
        
        \vspace{0.4cm}
        \item \textbf{ABAC} is highly \textbf{dynamic and flexible}, allowing for \textbf{fine-grained} access control without creating countless specific roles.
        
        \vspace{0.2cm}
        \begin{itemize}
            \large
            \item \textbf{Example Policy}: \\
            \textit{"Only managers \textbf{(Subject)} can read \textbf{(Action)} financial reports \textbf{(Object)} during business hours \textbf{(Environment)}."}
        \end{itemize}
    \end{itemize}
\end{frame}
% --- Slide: Attribute-Based Access Control (3) - Architecture ---
\begin{frame}{Attribute-Based Access Control (ABAC) (3)}
    \Large
    \textbf{ABAC Architecture (How it works)}
    \vspace{0.3cm}

    \begin{center}
    \resizebox{0.9\textwidth}{!}{
    \begin{tikzpicture}[
        % Định nghĩa style cho các khối
        box/.style={rectangle, draw=black, thick, fill=cyan!10, minimum width=2.5cm, minimum height=1cm, align=center, font=\bfseries},
        database/.style={cylinder, draw=black, thick, aspect=0.25, fill=gray!10, minimum width=2cm, minimum height=1.5cm, align=center, shape border rotate=90, font=\footnotesize},
        arrow/.style={->, >=stealth, ultra thick}
    ]
        % --- CÁC NODE (KHỐI) ---
        \node[box, fill=gray!20, rounded corners] (user) {User};
        \node[box, right=2cm of user] (pep) {PEP\\{\small (Enforcement)}};
        \node[box, right=2cm of pep, fill=red!10, rounded corners] (resource) {Resource};
        
        \node[box, below=1.5cm of pep, fill=yellow!20] (pdp) {PDP\\{\small (Decision)}};
        
        \node[database, left=1cm of pdp] (pip) {Attributes\\(Subject, Obj, Env)};
        \node[database, right=1cm of pdp] (pap) {Policies\\(Rules)};

        % --- CÁC MŨI TÊN ---
        \draw[arrow] (user) -- node[above, font=\small] {Request} (pep);
        \draw[arrow] (pep) -- node[above, font=\small] {Allow/Deny} (resource);
        
        \draw[arrow, <->] (pep) -- node[right, font=\small] {Evaluate} (pdp);
        
        \draw[arrow] (pip) -- (pdp);
        \draw[arrow] (pap) -- (pdp);
    \end{tikzpicture}
    }
    \end{center}

    \vspace{0.4cm}
    \begin{itemize}
        \normalsize
        \setlength{\itemsep}{0.1cm}
        \item \textbf{PEP (Policy Enforcement Point)}: Intercepts the access request.
        \item \textbf{PDP (Policy Decision Point)}: The "brain" that checks Attributes against Policies to make a \textit{Yes/No} decision.
    \end{itemize}
\end{frame}

\section[Context-Based Access Control]{Context-Based \\ Access Control (CBAC)}


% --- Slide 1: Concept & Core Idea ---
\begin{frame}{Context-Based Access Control (CBAC) (1)}
    \Large
    \textbf{Core Concept}
    \vspace{0.3cm}
    \begin{itemize}
        \item \textbf{CBAC} is an access control model that regulates access based on the \textbf{environmental context} or state of the system at the time of the request.
        \item While RBAC asks "Who are you?" and ABAC asks "What are your attributes?", CBAC emphasizes the \textbf{"Situation"} (Context).
        \item \textbf{Principle:} A user might have valid permissions, but if the request occurs in an abnormal context (e.g., wrong time, strange location), access is denied.
    \end{itemize}
\end{frame}

% --- Slide 2: Contextual Factors ---
\begin{frame}{Context-Based Access Control (CBAC) (2)}
    \Large
    \textbf{Key Contextual Factors}
    \vspace{0.3cm}
    \begin{itemize}
        \item \textbf{Time:} Allowing access only during business hours or specific working shifts.
        \item \textbf{Location:} Checking internal IP addresses, GPS coordinates, or geographical boundaries (Geofencing).
        \item \textbf{Device:} Distinguishing company-issued devices (Managed) from personal devices (BYOD), or checking security patch levels.
        \item \textbf{Network:} Requiring a secure VPN connection or access from a wired internal network.
    \end{itemize}
\end{frame}

% --- Slide 3: Comparison ---
\begin{frame}{Context-Based Access Control (CBAC) (3)}
    \Large
    \textbf{Comparing Access Control Models}
    \vspace{0.4cm}
    \begin{table}[]
        \centering
        \resizebox{\textwidth}{!}{
        \begin{tabular}{|l|l|l|}
        \hline
        \rowcolor{gray!20} \textbf{Model} & \textbf{Evaluation Focus} & \textbf{Characteristics} \\ \hline
        \textbf{RBAC} & Roles & Static, easy to manage but lacks flexibility. \\ \hline
        \textbf{ABAC} & Attributes & Highly granular but can be complex to implement. \\ \hline
        \textbf{CBAC} & \textbf{Environment (Context)} & \textbf{Dynamic, enhances real-time security.} \\ \hline
        \end{tabular}
        }
    \end{table}
    \vspace{0.3cm}
    \normalsize
    \textit{*Practice: CBAC is often used as an additional security layer on top of RBAC or ABAC.*}
\end{frame}

% --- Slide 4: Benefits & Scenario ---
\begin{frame}{Context-Based Access Control (CBAC) (4)}
    \Large
    \textbf{Benefits and Practical Scenario}
    \vspace{0.3cm}
    \begin{itemize}
        \item \textbf{Benefit:} Mitigates Account Takeover attacks. Even with a stolen password, hackers are blocked if accessing from an unusual country or unrecognized device.
        \item \textbf{Banking Scenario:}
        \begin{itemize}
            \large
            \item A loan officer can only view customer records when at the \textbf{branch} and during \textbf{business hours}.
            \item If accessed from \textbf{home} at \textbf{2 AM}, the system automatically locks the access.
        \end{itemize}
        \item \textbf{Trend:} Serves as a foundation for the modern \textbf{Zero Trust} security model.
    \end{itemize}
\end{frame}

\begin{frame}[plain] 
    \begin{center}
        \vspace{1.5cm}
        \Huge \textbf{Thank You!} \\
        \vspace{0.5cm}
        \Large for your listening and attention.
        
        \vspace{1.5cm}
        \LARGE \textbf{Q \& A} \\
        \vspace{0.3cm}
        \normalsize \textit{Do you have any questions for Group 6?}
    \end{center}
\end{frame}